% Generated draft: requires mathematical and source review.
\documentclass[11pt]{article}
\usepackage[margin=25mm]{geometry}
\usepackage{fontspec}
\setmainfont{Latin Modern Roman}
\usepackage{amsmath,amssymb,mathtools,unicode-math}
\setmathfont{Latin Modern Math}
\usepackage{enumitem,needspace,xurl,hyperref}
\setlength{\parindent}{0pt}
\setlength{\parskip}{5pt}
\setlength{\emergencystretch}{3em}
\newcommand{\N}{\mathbb{N}}
\newcommand{\Z}{\mathbb{Z}}
\newcommand{\Q}{\mathbb{Q}}
\newcommand{\R}{\mathbb{R}}
\newcommand{\C}{\mathbb{C}}
\newcommand{\F}{\mathbb{F}}
\newcommand{\A}{\mathbb{A}}
\newcommand{\PP}{\mathbb{P}}
\newcommand{\E}{\mathbb{E}}
\newcommand{\eps}{\varepsilon}
\newcommand{\dd}{\,\mathrm{d}}
\begin{document}
\section*{UnsolvedMath ID: 30003449}
\textbf{Source problem code:} catalogue entry
\subsection{Definitions and mathematical statement}
Under the standard complex-Hermitian convention, let \(g,d\geq 1\) be integers and let \(A_1,\ldots,A_g\in M_d(\C)\) satisfy \(A_j^*=A_j\). A Hermitian monic linear pencil is
\[
L(x)=I_d+\sum_{j=1}^g A_jx_j.
\]
For \(n\geq 1\), write \(\mathcal H_n=\{T\in M_n(\C):T^*=T\}\). For a tuple \(X=(X_1,\ldots,X_g)\in\mathcal H_n^g\), define
\[
L(X)=I_d\otimes I_n+\sum_{j=1}^g A_j\otimes X_j,
\qquad
D_L(n)=\{X\in\mathcal H_n^g:L(X)\succeq0\}.
\]
The free spectrahedron is the graded set \(D(L)=(D_L(n))_{n\geq1}\). Taking boundary level by level, let \(\partial D_L(n)\) denote the topological boundary in the real vector space \(\mathcal H_n^g\), and set
\[
\partial^1D(L)=\bigsqcup_{n\geq1}
\{X\in\partial D_L(n):\dim_{\C}\ker L(X)=1\}.
\]
Here the kernel is that of the evaluated \(dn\times dn\) matrix; the superscript \(1\) refers to kernel dimension, not matrix level.

\textbf{Extracted question.} For every Hermitian monic linear pencil \(L\) covered by the intended source hypotheses, must there exist an integer \(n\geq1\) and a tuple \(X\in\partial D_L(n)\) such that
\[
\dim_{\C}\ker L(X)=1?
\]
Equivalently, must \(\partial^1D(L)\neq\varnothing\)?

\textbf{Unresolved qualifications.} The definitions above explicitly adopt complex-Hermitian tuples and levelwise topological boundary. The supplied record does not specify the scalar field, the boundary convention, or whether the matrix level is unrestricted. A real-symmetric interpretation would instead use symmetric real matrices and real kernel dimension. The record states no minimality, irreducibility, nonconstancy, boundedness, or other nondegeneracy hypothesis on \(L\). In particular, the question is expressed in terms of the chosen pencil, and the record does not say whether redundant or repeated blocks are excluded. These potentially material restrictions must be checked before treating the unrestricted formulation as the intended conjecture; none is imposed here as an established source hypothesis.
\subsection{Short English statement}
The question asks whether a free spectrahedron has, at some matrix size, a boundary tuple at which its defining pencil has a one-dimensional kernel. The supplied record labels the question open, but its intended hypotheses and matrix-domain conventions require clarification. This is a statement-only draft requiring human mathematical review.
\subsection{Sources}
\begin{itemize}
\item Supplied local UnsolvedMath registry record, permanent ID \(30003449\), title \emph{Rank-One Kernel Boundary Points of Free Spectrahedra}, collection \emph{Oberwolfach Reports Open Problems}, exact source code \texttt{catalogue entry}. The record gives proposed year \(2017\) and names no individual proposer.
\item Supplied source citation: \emph{Real Algebraic Geometry With a View Toward Moment Problems and Optimization} (2018), Oberwolfach Reports, pp.\ 771--862, DOI \texttt{10.4171/owr/2017/14}. Original Oberwolfach report: \url{https://doi.org/10.4171/owr/2017/14}. The supplied background describes these reports as documenting mathematical workshops and their open-problem discussions.
\item Supplied literature source, labeled ``Laffey counterexample to Kippenhahn's conjecture'': \url{https://doi.org/10.1016/0024-3795(83)90159-3}.
\item Supplied literature source, labeled ``Free loci of matrix pencils'': \url{https://arxiv.org/abs/1512.02648}.
\item Supplied dated literature curation, checked \texttt{2026-08-21}: status ``open.'' Its assessment says that this boundary question is related to Kippenhahn's conjecture without asserting an equivalence; it attributes a disproof of the classical Kippenhahn conjecture to Laffey and says that this does not thereby settle the extracted boundary question. This assessment is supplied dated curation derived from the OWR record, not independently verified here; its time-sensitive status and cited claims require independent verification.
\item The registry labels the statement \texttt{corrected\_verified} and says that mathematical content was preserved while notation, typography, and sentence structure were repaired for publication. These are supplied editorial claims, not an independent mathematical verification of this draft or its hypotheses.
\end{itemize}
\end{document}
